\documentclass[11pt]{article}

\usepackage[margin=1in]{geometry}
\usepackage{amsmath,amssymb,amsthm,mathtools}
\numberwithin{equation}{section}
\usepackage{booktabs,tabularx,array}
\usepackage{graphicx}
\usepackage{caption}
\usepackage{float}
\usepackage{enumitem}
\usepackage{xcolor}
\usepackage{tikz}
\usetikzlibrary{positioning,arrows.meta,shapes.geometric}
\usepackage[round]{natbib}
\usepackage{microtype}
\usepackage{setspace}
\onehalfspacing
\usepackage{titlesec}

% Compact run-in headings within each numbered subsection.
\titleformat{\subsubsection}[runin]
  {\normalfont\normalsize\bfseries}{\thesubsubsection}{0.6em}{}
\titlespacing*{\subsubsection}{0pt}{0.6\baselineskip}{0.6em}
\usepackage{hyperref}

\graphicspath{{../figures/}{figures/}}
\hypersetup{
  colorlinks=true,
  linkcolor=blue!50!black,
  citecolor=blue!50!black,
  urlcolor=blue!50!black
}

\setlength{\parindent}{0em}
\setlength{\parskip}{0.45em}
\setlength{\emergencystretch}{3em}
\setlist[itemize]{leftmargin=1.3em,itemsep=0.12em,topsep=0.12em}
\setlist[enumerate]{leftmargin=1.45em,itemsep=0.12em,topsep=0.12em}
\allowdisplaybreaks

\newtheorem{proposition}{Proposition}
\newtheorem{lemma}{Lemma}
\newtheorem{corollary}{Corollary}
\theoremstyle{definition}
\newtheorem{definition}{Definition}
\newtheorem{assumption}{Assumption}
\newtheorem{remark}{Remark}

\newcommand{\dd}{\mathrm d}
\newcommand{\E}{\mathbb E}
\newcommand{\1}{\mathbf 1}

\title{Declining Fertility and Housing Costs}
\author{Tommaso De Santo\thanks{New York University (td2248@nyu.edu). I am grateful to Corina Boar, Guido Menzio, Raquel Fernández and Sharon Traiberman for their guidance and support. I thank David Berger, Jaroslav Borovička, Nicoló Ceneri, Gian Luca Clementi, Simon Gilchrist, Paolo Martellini, Virgiliu
		Midrigan, Abdoulaye Ndiaye, Alessandra Peter, Diego Perez, Martin Rotemberg, Saverio Spinella, Venky Venkateswaran, and participants to the NYU Macroeconomics Lunch Seminar and NYU Student Macroeconomics Lunch Seminar for their comments and suggestions.}}
\date{October 19, 2026\\[0.4em]\small\href{https://tommasodesanto.github.io/jmp.pdf}{Click here for the latest version}}

\begin{document}

\maketitle
\begin{abstract}
\noindent
I study how the cost of family space shapes fertility and what housing
policy can do about it. Housing costs are at the center of the public debate
on declining fertility: children need space, and young households, who have
little savings, find that space costly to pay for. I build a life-cycle model
in which households choose whether to have children, whether and how much to
rent or own, and how much to save. Buying requires a down payment, and
earnings are risky. I calibrate the model to the United States in 2007 and
treat the fall in births since then as an exogenous decline in the taste for
children. Along the transition to a new steady state with a smaller
population, fewer households make housing cheaper, which feeds back into
fertility. Because a child is a long commitment, earnings risk weighs on
fertility, and lower risk raises births. I then use the model to assess
housing policy. Easier mortgage credit could give young households access to
the larger homes a child requires, but it does little for births. Policies
that change long-term incentives, such as more housing supply and credit for
parents, raise fertility, and their effects build up over a generation.
Housing policy must therefore be judged by its long-run effects.

\end{abstract}
\vspace*{\fill}
\pagebreak
\end{document}
